\chapter{Avleggaren}

\begin{motto}
Frenologien forsvann ikkje fordi han var vond.\\ Han forsvann fordi nokon bygde nevrovitskapen ved sida av.
\end{motto}

Det finst ein kategori i vitskapshistoria som designfaget høyrer heime i, og som det aldri har vedkjent seg. Det er kategorien av felt som ein gong var respekterte, som hadde utøvarar, institusjonar, litteratur og prestisje, og som forsvann; ikkje fordi dei vart motbeviste i ein einaste avgjerande test, men fordi det vaks fram eit fag med eit fundament ved sida av dei, og det nye faget gjorde det gamle overflødig. Frenologien hadde tidsskrift, foreiningar, framståande utøvarar; han forsvann då nevrovitskapen kom. Alkemien hadde tusen års akkumulert praksis; han forsvann då kjemien kom. Vitalismen forklarte livet med ei eigen livskraft; han forsvann då biokjemien viste at livet er kjemi heile vegen ned. Astrologien kartla himmelen med stor presisjon; astronomien overtok himmelen og lét astrologien sitje att med teikna.

Det desse faga deler er ikkje at dei var dumme. Mange kloke menneske arbeidde i dei, og dei produserte reelle observasjonar. Det dei deler er ein \emph{struktur}: dei hadde ein katalog utan ein teori, eit mønstergjenkjennande blikk utan eit fundament som kunne grunngje mønstera, og difor kunne dei ikkje skilje dei sanne mønstera frå dei innbilte. Frenologen såg verkeleg samanhengar mellom hovud og karakter; han hadde berre ingen teori som kunne seie kva for samanhengar som var reelle. Det er nett designfaget si stode. Det ser verkelege mønster i form; det manglar berre teorien som kan seie kvifor dei er der, og difor kan det ikkje skilje innsikt frå smak, kvalitet frå mote, kompetanse frå konvensjon. Denne boka har vist den same strukturen i kvar epoke av faget si historie.

\section{Eit fag som valde å ikkje bli paradigmatisk}

Thomas Kuhn gav oss ordet for tilstanden \autocite{kuhn1962}. Eit \emph{før-paradigmatisk} fag er eit fag før det har funne sitt fyrste delte rammeverk: konkurrerande skular som ikkje kan avgjere usemjene sine, kvar generasjon som byrjar grunnlaget på nytt, ingen delt mengd problem, ingen akkumulasjon. Kvar setning i denne skildringa passar på designfaget, og dei føregåande kapitla har dokumentert kvar einskild: skular frå Bauhaus til Ulm til den spekulative til design thinking, som ikkje kan avgjere usemjene sine fordi det ikkje finst eit grunnlag å avgjere dei på.

Men her møter faget eit kjært forsvar, og det er verdt å ta på alvor. Forsvaret lyder: «Kuhn skreiv om naturvitskapane. Samfunns- og humanvitskapane er ikkje før-paradigmatiske; dei er \emph{multi}-paradigmatiske, og det er ein moden tilstand, ikkje ein umoden. Design høyrer til der.» Innvendinga byggjer på ei reell usemje i Kuhn-eksegesen. Somme lesarar held fast ved Kuhn sjølv: eit fag er anten før sitt fyrste paradigme eller etter det, og striden mellom skular er nett kjenneteiknet på den før-paradigmatiske tilstanden \autocite{kuhn1962}. Andre meiner det er for grovt for fag der fleire rammeverk lever side om side i fruktbar spenning, og kallar slike fag multi-paradigmatiske. La oss for argumentet sin skuld gje den andre lesinga rett. Hjelper det design?

Det gjer det ikkje. Eit multi-paradigmatisk fag har \emph{paradigme}, fleire av dei, kvar med sitt presise rammeverk, sine omgrep, sine kriterium for kva som tel som eit godt arbeid innanfor det paradigmet. Ein økonom veit kva ein neoklassisk modell krev og kva ein åtferdsøkonomisk modell krev; han kan halde båe, veksle mellom dei, og vurdere eit arbeid \emph{innanfor} kvart. Det finst grunn å stå på, fleire grunnar, kvar fast. Det design har, er ikkje fleire fundament. Det er \emph{ingen}. Skulane i designhistoria er ikkje rivaliserande paradigme med kvar sitt presise grunnlag; dei er, som boka har vist kapittel for kapittel, rivaliserande \emph{haldningar}, kvar utan eit fundament som kunne gjere usemja avgjerbar innanfrå. Eit multi-paradigmatisk fag kan seie kvifor eit arbeid er godt etter eit av paradigma sine. Design kan ikkje seie det etter noko, fordi det ikkje har eit «etter» å seie det etter.

Så design er verken det eine eller det andre i vanleg forstand. Det er ikkje før-paradigmatisk på veg mot eit paradigme, slik kjemien var før Lavoisier. Det er ikkje multi-paradigmatisk med fleire faste grunnar, slik økonomien er. Det er ein tredje ting, og den verste: eit fag som har \emph{vald} å ikkje bli paradigmatisk. Det har teke den før-paradigmatiske striden mellom skular og gjort han permanent, ikkje ved å løyse han og ikkje ved å institusjonalisere fleire grunnar, men ved å omdøype sjølve grunnløysa til ein dygd. «Design er ein eigen kunnskapsform.» «Problema våre er vrange.» Dette er ikkje teoriar om eit fag i rørsle mot eit paradigme; det er erklæringar om at faget aldri har tenkt å reise eitt, kledde som metodologiske innsikter. Herbert Simon baud fram eit fundament, ein vitskap om det kunstige med presise omgrep om søk og avgrensa rasjonalitet \autocite{simon1969sciences}, og faget tok imot prestisjen i tittelen og lét disiplinen liggje. Eit ungt fag som ikkje har funne paradigmet sitt, fortener tolmod. Designfaget er ikkje ungt. Det har vore akademisk i over hundre år, og har hatt all den tid eit fag treng for å finne fundamentet sitt. Eit gamalt fag som har funne unnskyldningar for å sleppe å leite, fortener diagnosen: ikkje multi-paradigmatisk, men a-paradigmatisk av eige val. Det er ein avleggar-tilstand, ikkje ein modnings-tilstand.

\section{Fingrane på hovudet}

Eg har sagt at frenologien hadde tidsskrift og foreiningar, og at det er der parallellen ligg. Så lenge det står som ein laus påstand, kan ein design\-akademikar trygt avfeie det som retorikk. Han kan ikkje avfeie tala, og han kan ikkje avfeie andletet.

Franz Joseph Gall (1758--1828) grunnla det han kalla organologi, læra om at hjernen er sett saman av organ som kvar svarar til ein eigenskap, og at storleiken på organet kan lesast av forma på skallen utanpå. Eleven hans, Johann Spurzheim, gav læra namnet «frenologi», skilde lag med meisteren og bar henne til Storbritannia og USA som ein omreisande apostel. Då George Combe i 1828 gav ut \emph{The Constitution of Man}, fekk rørsla si folkebok: ho vart ein av dei mest seldte titlane i heile hundreåret, i opplag som måler seg med Bibelen i somme marknader, og ho gav frenologien noko meir enn tilhengjarar. Ho gav han \emph{autoritet}, den slags brei kulturell tiltru eit felt berre får når det er gått inn i kvardagsspråket. Folk tala om kvarandre sine «organ» slik vi i dag talar om kvarandre sin smak.

Og rørsla bygde institusjonar, presist dei eit fag «skal» ha. Edinburgh Phrenological Society, grunnlagt i 1820, var det fyrste frenologiske selskapet i verda; etter det kom selskap i by etter by over heile den engelsktalande verda. Frå 1823 gav Edinburgh-miljøet ut \emph{The Phrenological Journal and Miscellany}, eit fagtidsskrift som heldt fram i tiår, med artiklar, ordskifte, fagfellar som las kvarandre. I New York dreiv firmaet Fowler \& Wells frå 1838 \emph{The American Phrenological Journal}, som gjekk heilt ut på 1900-talet, eit lengre liv enn mange respekterte vitskaplege tidsskrift har hatt. Fowlerane gjorde meir enn å skrive: dei opna ein praksis. For ein avgift kjende dei på hovudet ditt og las karakteren din av buklane, «practical phrenology», ein betalt klinisk konsultasjon med diagnose og råd, akkurat som ein lækjar eller ein rådgjevar. Grunnleggjar og lærelinje. Populærbok med massiv autoritet. Lærde selskap. Fagfelle-tidsskrift som varer i generasjonar. Betalt yrkespraksis med klientar. Internasjonal spreiing. Set lista ved sida av det føregåande kapitlet sin gjennomgang av designfaget sitt apparat, akkrediteringa, prisane, biennalane, tidsskrifta, professorata, og dei to listene er den same lista. Alt det ytre apparatet eit fag kan ha, hadde frenologien, og likevel forsvann han. Apparatet er ikkje fundamentet.

Gjev det heile eit andlet, for ein avleggar er ikkje ein statistikk; han er ein mann som trur han les noko verkeleg, og som har rett i at han ser eit mønster og feil i alt han trur mønsteret tyder. Set deg i Edinburgh ein kveld kring 1830, i ein salong med tunge gardiner og levande ljos, der dei vitelystne i byen har samla seg. Ein frenolog, kanskje eit medlem av Edinburgh Phrenological Society, står bak ein stol der ein gjest har sett seg. Han ber gjesten løyse opp håret, og så byrjar det: fingrane går over skallen, sakte, metodisk, frå panna bakover, kjenner etter buklar og søkk, dvelar ved kvar oppløfting. Han har eit kart i hovudet, ein byste av porselen med felt teikna opp, kvart merkt med ein eigenskap: Vørdnad her, Stridslyst der, Velvilje, Vinningssyke, Von. Der handa kjenner ei oppløfting, er organet under stort, og eigenskapen sterk.

Sjå kor overtydande det er. «De har eit stort organ for Velvilje», seier frenologen med fingrane framleis på hovudet, «De er ein gjevmild natur.» Gjesten kjenner seg sett, for kven vil ikkje vere gjevmild, og nikkar; salongen nikkar med. «Stridslyst er måteleg utvikla, men Fasthald er sterkt: De gjev ikkje opp lett.» Igjen treff det, fordi skildringa er romsleg nok til å femne kven som helst, fordi gjesten leitar etter stadfesting og finn ho, fordi salongen vil at det skal stemme. Og når noko \emph{ikkje} stemmer, når gjesten protesterer at han slett ikkje er sparsam, har frenologen alltid ei forklaring: organet er der, men eit anna held det i sjakk, eller oppsedinga har dempa det medfødde. Inga lesing kan bomme. Det er ikkje juks; mannen trur. Han har eit kart, ei lære, eit selskap, eit tidsskrift, og no, under fingrane, ein skalle som ser ut til å svare. Han har alt utan det eine som ville late han skilje det han kjenner frå det han innbiller seg: ein teori som kunne seie kva for buklar som tyder noko og kva for nokre som berre er bein.

Spol så fram nokre tiår, til det rommet der alt han gjorde vart fortrengt. Det skjer ikkje i ein salong, men på eit obduksjonsbord i Paris. Då Paul Broca i 1861 bøygde seg over hjernen til ein avliden mann som i levande live hadde mist evna til å tale, kunne forstå alt men ikkje få fram meir enn ei einaste staving, fann han ein skade, ein mjukna flekk, i ein bestemt vinding av venstre pannelapp. Han hadde ikkje kjent på ein skalle utanpå; han hadde opna han, og funne korleis ein funksjon, talen, heng saman med ein bestemt stad i vevet, prova ved at staden var øydelagd der funksjonen var borte. Carl Wernicke fann i 1874 eit anna språkområde på liknande vis. For fyrste gong hadde nokon lokalisert funksjon i hjernen \emph{empirisk og teoretisk fundert}, ikkje ved å lese forma utanpå, men ved å knyte skade til tap inne i vevet, med ein teori som kunne seie kvifor sambandet var reelt. Frenologen hadde sagt at funksjon sit på bestemte stader; det var ikkje galt. Men han las stadene av buklar utanpå skallen, utan ein måte å vite om han las rett. Nevrologen las dei av sjølve vevet, med ein måte å vite. Same prosjekt, kart over hjernen; den eine utan fundament, den andre med.

Og legg merke til kva som \emph{ikkje} skjedde den dagen på obduksjonsbordet. Frenologien fall ikkje død om. Selskapa heldt fram med å møtast; tidsskrifta kom framleis ut; det fanst framleis frenologar med fingrane på hovud i salongar i fleire tiår etter. Avleggaren veit ikkje straks at han er avleggar. Han veit det ikkje før verda sakte sluttar å spørje han, fordi nokon ved sida av no svarar betre. Dette er designfaget sin situasjon, ord for ord: det har apparatet, like fullstendig som frenologien sitt, og det manglar fundamentet, like fullstendig. Det arbeider, som dei føregåande kapitla har vist, like hardt som frenologen med å bortforklare kvar bom og rekne kvart treff som stadfesting. Det einaste spørsmålet som står att, er det same som for frenologien: kva tid noko betre står opp ved sida av.

\begin{kasus}{Lesinga som alltid stemte}
Sett frenologen si lesing opp som ein test. Han kjenner ein bukl over øyret og seier «sterk Stridslyst». Tre utfall er moglege. Gjesten er stridig: treff, læra stadfest. Gjesten er mild: men sjå, eit stort organ for Velvilje held Stridslysta i sjakk, læra stadfest ved eit anna organ. Gjesten har ingen merkbar legning: medfødd anlegg, dempa av oppseding, læra stadfest ved ein tredje mekanisme. Det finst ikkje noko utfall som tel mot læra. Ein påstand som ingen observasjon kan svekke, måler ingenting; han berre kler kvar observasjon i sine eigne ord. Det er ikkje at frenologen lyg om kva han kjenner; han kjenner verkeleg buklar, og det \emph{finst} samband mellom hovudform og eitt og hitt. Det er at han ikkje har nokon måte å skilje sambandet som er reelt frå mønsteret han legg på. No byt ut frenologen med designaren og buklen med ei form. Designaren ser eit arbeid og seier «dette sit, dette har integritet». Arbeidet sel: treff, dommen stadfest. Arbeidet sel ikkje: men det var forut for si tid, dommen stadfest. Arbeidet gjer verken frå eller til: smaken vil kome, dommen stadfest. Same struktur, ord for ord, og fråvêret av den eine tingen som ville gjort det til kunnskap. Frenologen fekk høyre at han var avleggar då nokon opna vevet. Spørsmålet til designaren er kva tid nokon opnar forma.
\end{kasus}

\section{Materialet som svarar tilbake}

Avleggarane forsvann ikkje då dei vart motbeviste; dei forsvann då verda fekk bruk for noko dei ikkje kunne levere, og eit anna fag kunne. Den krafta er no på veg mot designfaget, og ho kjem frå materialet sjølv. Materialet finst alt, og det ber tesen tydelegare enn noko argument.

I heile faget si historie har formgjevaren arbeidd med passivt stoff: tre, metall, plast, pikslar, materiale som gjorde nett det dei vart sagt. No endrar materialet seg under hendene våre. I 2020 sette Michael Levin og Josh Bongard, ved Tufts og University of Vermont, saman det dei kalla «xenobots»: levande maskiner bygde av hudceller og hjarteceller frå froskearten \emph{Xenopus laevis}, sette saman ikkje av ein formgjevar med eit blikk, men av ein evolusjonær algoritme som søkte gjennom rommet av moglege kroppsformer etter ein som kunne røre seg mot eit mål. Cellene, frigjorde frå froskekroppen dei skulle ha vorte til, organiserte seg etter den fora algoritmen fann, og gjorde noko froskeceller aldri «skal» gjere: dei symde, dei skuva småpartiklar i samla rørsle, dei samarbeidde. I 2021 viste same gruppa noko endå meir urovekkjande, kinematisk sjølvreplikasjon: xenobotane gjekk omkring i ein skål med lause celler, sopa dei saman i haugar, og haugane vart til nye xenobotar, som i sin tur kunne sope. Ei form som lagar fleire av seg sjølv, utan at nokon teikna den neste.

Spør deg kva designfaget sitt einaste reiskap, det mønstergjenkjennande blikket, kan gjere her. Ingenting. Du kan ikkje sjå på ein klump froskeceller og intuere deg fram til kva form som vil symje, langt mindre kva form som vil \emph{replikere}. Det rommet er for stort, og det er ikkje styrt av smak; det er styrt av kva cellene, med sin eigen agenda, faktisk gjer under eit gjeve seleksjonstrykk. For å arbeide her treng du ein modell av kva ein agent er, korleis kompetanse og mål fordeler seg nedover skalaene frå organisme til celle, kva eit formrom er og korleis ein søkjer i det \autocite{levin2022,fields2022}. Du treng eit fundament. Det atelier-blikket som var faget sin stoltheit, det som skulle skilje den gode forma frå den dårlege ved å sjå, er her ikkje berre utilstrekkeleg; det er irrelevant. Materialet svarar tilbake, og blikket har ingenting å seie til eit materiale som har sin eigen vilje.

Og medan materialet vert agentisk i den eine enden, vert sjølve formframbringinga automatisert i den andre. Når generative system kan produsere formene, tusenvis av variantar i timen, kvar kompetent på handverkets premiss, då er handverket, faget sin siste eigentlege kompetanse, ikkje lenger faget sitt einaste. Då sit faget att med berre det det aldri bygde: teorien om kvifor \emph{denne} forma er betre enn dei tusen andre maskina nett la fram. Det er den einaste oppgåva som står att når både handa og no materialet er teke ut av faget sine hender, og det er nett den oppgåva faget i hundre år har valt å ikkje løyse. Frenologen heldt det gåande så lenge alt han hadde, var daude skallar å kjenne på. Gjev han ei levande celle som navigerer, og handa hans har ingenting å lese. Designfaget har fått si celle.

Lat meg vere presis om kva eg spår, for eg spår ikkje at design forsvinn. Verda vil alltid trenge at nokon gjev form til ting; det er like umisteleg som at verda treng at nokon lækjer sjuke. Det eg spår er at \emph{faget slik det er}, den før-paradigmatiske disiplinen som deler ut grader i ein kompetanse han ikkje kan definere, kjem til å bli avleggar, slik medisinen før fysiologien vart avleggar medan lækjinga heldt fram. Formgjevinga overlever. Spørsmålet er kva fag som kjem til å eige henne, og her er det to vegar, og berre to. Anten finn formgjevinga sitt fundament innanfrå: designfaget gjer endeleg den overgangen det har utsett i hundre år, byggjer eller adopterer ein presis teori om form, og vert eit paradigmatisk fag som fortener plassen sin. Eller så vert formgjevinga overteken utanfrå, av informatikken, av materialvitskapen, av eit komande felt for agentisk formgjeving som har det fundamentet designfaget mangla, og designfaget sit att med teikna, slik astrologen sat att med stjerneteikna då astronomen overtok himmelen. Vi har alt sett poda flytte ut éin gong: design thinking migrerte til økonomifaget fordi det ikkje fanst grunn å bli verande heime. Det same kan skje med heile faget. Begge vegar endar den noverande disiplinen; den eine let han bli til noko betre, den andre til ein fotnote. Det faget ikkje lenger har, er den tredje vegen: å halde fram som før.

\vignett

Klagemålet er ført, og historia er fortald. Det står att eitt spørsmål, og berre eitt: finst det ein grunn å byggje på, om faget vil? Det gjer det. Det siste kapitlet peikar mot han.

\laeringsmal{\begin{itemize}
\item Gjere greie for kategorien «avleggar»: eit fag som ikkje fell for eit motbevis, men fordi eit fag med fundament veks opp ved sida av, og plassere frenologi, alkemi, vitalisme og astrologi i han.
\item Vise, med frenologien sin institusjonshistorie (Gall, Spurzheim, Combe, dei frenologiske selskapa og tidsskrifta, Fowlerane sin praksis) som døme, kvifor eit fullt ytre apparat ikkje er det same som eit fundament.
\item Skilje mellom før-paradigmatisk, multi-paradigmatisk og a-paradigmatisk, og grunngje kvifor design høyrer til den siste: eit fag som har \emph{vald} å ikkje bli paradigmatisk.
\item Forklare kvifor agentisk materiale (xenobotane til Levin og Bongard) og generativ automatisering saman avslører fråvêret av ein teori om form, og kva for to vegar som då står att for faget.
\end{itemize}}

\ovingar{\begin{enumerate}
\item Sett opp to spalter: frenologien sitt apparat (selskap, tidsskrift, populærbok, betalt praksis) mot designfaget sitt apparat frå førre kapittel (akkreditering, prisar, biennalar, tidsskrift). Argumenter for eller imot at listene er den same lista, og avgjer kva ein slik likskap kan og ikkje kan bevise.
\item Frenologien fall ikkje for éin avgjerande test, men då Broca (1861) og Wernicke (1874) bygde lokaliseringa på eit fundament ved sida av. Finn eit anna avleggar-fag og avgjer kva som stod opp «ved sida av» det, og om mønsteret held.
\item Vurder designfaget sitt forsvar om at det er multi-paradigmatisk snarare enn før-paradigmatisk. Kva ville krevjast for at forsvaret skulle halde, og er kravet oppfylt?
\item Vel ei formgjevingsoppgåve og prøv fyrst å løyse henne med atelier-blikket, så å skissere korleis ho ville løysast med eit søk over eit formrom (slik xenobotane vart funne). Kva gjer blikket som søket ikkje gjer, og omvendt, og kva sit faget att med når materialet har sin eigen agenda?
\end{enumerate}}

\vidarelesing{\fullcite{kuhn1962}; \fullcite{simon1969sciences}; \fullcite{levin2022}; \fullcite{fields2022}; \fullcite{michl2014}.}
